\subsection{Exponential and logarithm}

We denote by exp the complex exponential function from C to C (FVR, III, p.~8, def. 2). It is differentiable and satisfies \(\exp' = \exp\) (FVR, III, p.~9, (26)); therefore it is holomorphic in C. Let A be a unital Banach algebra and \(a\) an element of A. By prop. 2 of I, p.~65 and formula (9) of FVR, III, p.~16, we have

\[
\exp (a) = \sum_ {n = 0} ^ {\infty} \frac {a ^ {n}}{n !}.\tag{9}
\]

Since \(\|a^{n}\|\leqslant\|a\|^{n}\) we see that \(\|\exp(a)\|\leqslant\exp(\|a\|)\) and that the series (9) converges uniformly in every ball of A. The map \(a\mapsto\exp(a)\) from A to A is holomorphic (prop. 10 of I, p.~76). One also sometimes denotes \(e^{a}\) the exponential of \(a\in A\) .

When \(a\) is an endomorphism of a Banach space E, the exponential \(\exp(a)\) thus defined in the Banach algebra \(\mathcal{L}(\mathrm{E})\) coincides with that defined in FVR, IV, p.~27, def. 1, according to loc. cit. prop. 7 (3).

For every element b of A that is permutable with a, we also have

\[
\exp (a + b) = \sum_ {n = 0} ^ {\infty} \frac {b ^ {n}}{n !} \exp (a),
\]

(prop. 11 of I, p.~77), whence

\[
\exp (a + b) = \exp (a) \cdot \exp (b).\tag{10}
\]

In particular, \(\exp(a)\) is invertible and

\[
\exp (a) ^ {- 1} = \exp (- a).\tag{11}
\]

Let B be the set of \(z \in C\) such that \(-\pi < Jz < \pi\) . Let F be the complement in C of the interval \(R_{-}\). The restriction of the exponential to B induces, by passing to subspaces, a bijection from B onto F (FVR, III, p.~10, n°7), whose inverse bijection will be denoted by log.

If \(a \in A\) is such that \(\mathrm{Sp}_{\mathrm{A}}(a) \subset \mathrm{F}\) , one can form the element \(\log(a)\) of A. We have \(\mathrm{Sp}_{\mathrm{A}}(\log(a)) \subset \mathrm{B}\) , and

\[
\exp (\log (a)) = a\tag{12}
\]

by prop. 8 of I, p.~75. Conversely, let \(b\) be an element of A such that \(\mathrm{Sp}_{\mathrm{A}}(b) \subset \mathrm{B}\) . We have \(\mathrm{Sp}_{\mathrm{A}}(\exp(b)) \subset \mathrm{F}\) and

\[
\log (\exp (b)) = b\tag{13}
\]

(loc. cit.).

In particular, if \(a \in A\) is such that \(\varrho(a) < 1\) , one has \(\mathrm{Sp}_{\mathrm{A}}(1 - a) \subset \mathrm{F}\) and we can form \(\log(1 - a)\) . For \(n \geqslant 1\) , the \(n\) th derivative of \(z \mapsto \log(1 - z)\) is \(z \mapsto -(n - 1)! (1 - z)^{-n}\). The power series expansion of \(z \mapsto \log(1 - z)\) at the point 0 is therefore

\[
\log (1 - z) = - \sum_ {n = 1} ^ {\infty} \frac {z ^ {n}}{n},
\]

valid for \(|z| < 1\) (VAR, R1, p.~30, 3.3.9). By prop. 2 of I, p.~65, we obtain

\[
\log (1 - a) = - \sum_ {n = 1} ^ {\infty} \frac {a ^ {n}}{n}.\tag{14}
\]

PROPOSITION 12. --- Let A be a commutative unital Banach algebra. The image of the exponential map is the connected component of the identity in the group G of invertible elements of A.

Formulas (10) and (11) prove that \(\exp(A)\) is a subgroup of G. By the preceding (see formula (12)), this subgroup contains the open ball of center 1 and radius 1. It is therefore an open subgroup, and consequently closed, of G. Moreover, A is connected and the map \(a \mapsto \exp(a)\) is continuous, so that \(\exp(A)\) is connected. Therefore \(\exp(A)\) is the identity component of G.

